\documentstyle[%


righttag%


,11pt%


,amssymb%


,russian%               remove in Eng.


,izvru%                 Set izven% in Eng.


]{amsart}





%       Math definitions





\newcommand{\tts}[1]{}





%\hoffset=-15mm%                  For Eng.


 


\setcounter{page}{26}


\god{1999}


\nomer{9~(448)} %                        Remove in Eng.


%\nomer{Vol.\,43, No.\,9}%               Set in Eng.


%\str{\pageref{begin}--\pageref{end}}%  Set in Eng.


\udk{513.8}


\author{\it А.Я.\,Нарманов}


% NB:   ^^ remove in Eng.


\title{О геометрии вполне геодезических римановых слоений}





\date{}


%\pagestyle{myheadings}%         Set in Eng.


\pagestyle{plain} %              Remove in Eng.


\begin{document}


%\thispagestyle{plain}%          Set in Eng.


\maketitle


\label{begin}








Пусть $M$ --- гладкое связное приводимое риманово многообразие. 


Тогда на $M$ существуют два параллельных слоения, взаимно 


дополнительных по ортогональности ([1], с.\,173). Если $M$ 


полно и односвязно, то имеет место теорема де~Рама, которая утверждает, 


что $M$ изометрично прямому произведению любых двух слоев из 


разных слоений ([1], с.\,180). В этом случае оба слоения являются 


римановыми и вполне геодезическими одновременно. В данной работе 


изучается риманово вполне геодезическое слоение на $M$,


причем полная интегрируемость дополнительного ортогонального


распределения не предполагается. 


В статье всюду гладкость означает гладкость класса $C^\infty$.


\medskip





{\bf 1.} Пусть $M$ --- гладкое связное полное риманово многообразие 


размерности $n$ с римановой метрикой $g$, $\nabla$ --- связность


Леви-Чивита, $F$ -- гладкое слоение размерности $k$ на $M$ ([2];


[3], c.\,24). 


Обозначим через $L(p)$ слой слоения $F$, проходящий через точку $p$, 


$F(p)$ --- касательное пространство слоя $L(p)$ в точке $p$, $H(p)$ ---


ортогональное дополнение $F(p)$ в $T_pM$, $p\in M$. Возникают два 


подрасслоения (гладкие распределения) $TF=\{F(p) : p\in M\}$, $H=\{H(p): 


p\in M\}$ касательного расслоения $TM$ такие, что $TM=TF\oplus H$, где 


$H$ является ортогональным дополнением $TF$. По определению 


слоения, для каждой точки $p\in M$ существует окрестность $U$ точки 


$p$ и локальная система координат


$(x^1, x^2, \dotsc, x^k, y^1, y^2, \dotsc, y^{n-k})$ 


на $U$ такие, что


$\{\frac {\partial}{\partial x^1},\frac {\partial}{\partial 


x^2}, \dotsc, \frac {\partial}{\partial x^k}\}$ образует базис


для гладких сечений 


$TF|_U$ (сужение $TF$ на $U$) [2]. Пусть $\omega^1, \omega^2, \dotsc, 


\omega^k$ --- гладкие дифференциальные $I$-формы, $\nu_1, \nu_2, \dotsc, 


\nu_m$ --- гладкие векторные поля на $U$ такие, что $\{\omega^i, 


dy^\alpha\}$ образуют локальный базис для сечений кокасательного 


расслоения, $\{\frac {\partial}{\partial x^i},\nu_\alpha\}$ образуют базис для 


сечений касательного расслоения $TU=TF|_U\oplus H|_U$, 


двойственной к $\{\omega^i, dy^\alpha\}$, где $m=n-k$,


$1\leq i\leq k$, $1\leq 


\alpha \leq m$. Такую окрестность обозначим через $U(x,y)$.





Теперь предположим, что $F$ является римановым слоением по 


отношению к $g$ [2], [3]. Это означает, что в каждой окрестности 


$U(x,y)$ риманова метрика $g$ имеет вид


$$


g_{i_j}(x, y)\omega^i \omega^j + g_{\alpha \beta}(y)dy^\alpha dy^\beta ,


$$


где $1 \leq i$, $j \leq k$, $1\leq \alpha$, $\beta \leq m$,


$x=(x^1, x^2, \dotsc, x^k)$, 


$y=(y^1, y^2, \dotsc, y^m)$.


\medskip





{\bf Замечание.} В целях упрощения обозначений, в выражениях, где имеется 


суммирование по повторяющимся индексам, будем опускать знак 


суммирования.


\medskip





Кусочно-гладкую кривую $\gamma : [0,1]\to M$ назовем горизонтальной, 


если $\Dot \gamma (t)\in H(\gamma (t))$ для каждого $t\in [0,1]$.


Кусочно-гладкая кривая, которая лежит в слое $F$, называется вертикальной.





Пусть $I = [0,1]$, $\nu :I\to M$ --- вертикальная кривая, $h:I\to M$ ---


горизонтальная кривая и $h(0)=\nu(0)$. 


Кусочно-гладкое отображение $P:I\times I\to M$ такое, что 


$t\to P(t,s)$ является 


вертикальной кривой для каждого $s\in I$, кривая $s\to P(t,s)$ является 


горизонтальной для каждого $t\in I$, причем $P(t,0)=\nu (t)$ для $t\in I$, 


$P(0,s)=h(s)$ для $s\in I$, называется вертикально-горизонтальной


гомотопией [4], [5]. Если для каждой пары вертикальной


и горизонтальной кривых $v, h:I\to M$ с $v(0)=h(0)$


существует соответствующая вертикально-горизонтальная 


гомотопия $P$, то говорят, что распределение 


$H$ является связностью Эресмана для слоения $F$ [5], [6]. 


Так как рассматриваемое слоение $F$ риманово, и многообразие 


$M$ полно, распределение $H$ является связностью 


Эресмана для $F$ [4], [5].


Для каждой 


окрестности $U=U(x,y)$ слоение $F|_U$ (сужение $F$ на $U$ ) 


задается гладкой субмерсией $f:U\to R^{n-k}$ [2]. Как показано в [5], 


сужение $H$ на $U$ является связностью Эресмана для 


сужения $F$ на $U$ тогда и только тогда, когда сужение $H$ на 


$U$ является связностью Эресмана для расслоения $f:U\to R^{n-k}$.





Пусть $\pi_1:TM\to TF$, $\pi_2:TM\to H$ --- ортогональные проекции,


$V(M)$, $V(F)$, $V(H)$ --- множество гладких сечений расслоений $TM$, 


$TF$, $H$ соответственно. Если $X\in V(F)$ $(X\in V(H))$, то $X$ 


назовем вертикальным (горизонтальным) полем.





Пусть каждый слой $F$ является вполне геодезическим 


подмногообразием $M$. Это эквивалентно тому, что $\nabla_XY\in V(F)$ 


для всех $X,Y\in V(F)$ ([3], с.\,47--61). В этом случае говорят, что $F$ 


является вполне геодезическим слоением. В дальнейшем будем 


предполагать, что $F$ является римановым слоением, слои которого 


являются вполне геодезическими подмногообразиями $M$. Тогда на 


расслоениях $TF$ и $H$ определены метрические связности 


$\nabla^1$ и $\nabla^2$ следующим образом. Если $X\in V(F)$, $Y\in V(H)$, 


$Z\in V(M)$, то


$$


\nabla_Z^1X=\pi_1(\nabla_ZX),\quad \nabla_Z^2Y=\pi_2[Z_1, \wt Y] + 


\pi_2[\nabla_{Z_2}Y],


$$


где $Z=Z_1\oplus Z_2$, $Z_1 \in V(F)$, $Z_2 \in V(H)$, $\wt Y \in V(M)$,


$\pi_2\wt Y=Y$;


здесь $[Z_1, \wt Y]$ --- скобка Ли векторных полей $Z_1$ и $\wt Y$.\;


$\nabla^2$ является метрической связностью тогда и только тогда, когда 


$F$ является римановым. В силу того, что $F$ вполне геодезично, 


$\nabla^1$ также является метрической связностью ([3], с.\,47--61;


[7], с.\,257).





Пусть $p\in M$, $S(p)$ --- множество точек $M$, которые можно соединять 


с $p$ горизонтальными кривыми. В силу того, что $F$ вполне 


геодезично, для каждого $p\in M$ множество $S(p)$ обладает топологией 


и дифференциальной структурой, по отношению к которым $S(p)$ 


является гладким погруженным подмногообразием $M$ ([6], теорема 4).





\begin{lem}


$\dim S(p)\ge n-k$ для каждой точки $p\in M$.


\end{lem}





\begin{pf} Пусть $p\in M$. Существует достаточно малое $\varepsilon 


>0$ такое, что $\exp_p:B_\varepsilon \to V$ является диффеоморфизмом, где 


$B_\varepsilon =\{X\in T_pM: |X|<\varepsilon \}$, $V$ --- открытая окрестность 


точки $p$, $|X|$ --- длина вектора $X$, $\exp_p$ --- экспоненциальное 


отображение в точке $p$. Положим $H_\varepsilon = B_\varepsilon \cap 


H(p)$, $Q(p)= \exp_p (H_\varepsilon )$. Так как $F$ риманово,


всякая геодезическая, ортогональная в некоторой своей точке


к слою $F$, остается ортогональной к слоям $F$ во всех своих 


точках ([2], с.\,123). Следовательно, 


$Q(p)\subset S(p)$.


\end{pf}





Пусть $P:I\times I\to M$ --- вертикально-горизонтальная гомотопия. 


Обозначим через $D_sP(t,s)$ касательный вектор кривой $s\to P(t,s)$ в 


точке $P(t,s)$, через $D_tP(t,s)$ касательный вектор кривой $t\to P(t,s)$ в 


точке $P(t,s)$.





\begin{lem}


Пусть $X(t,s)=D_sP(t,s)$, $Y(t,s)=D_tP(t,s)$, для $(t,s)\in I\times I$. 


Тогда $\nabla_X^1Y = 0$ и $\nabla^2_YX=0$.


\end{lem}





\begin{pf}


Нетрудно проверить, что $[X,Y]=0$. Так как $Y$ ---


вертикальное поле, имеем $\nabla^2_YX=\pi_2[X,Y]$. Известно, 


что$\nabla^1_ZY=\pi_1[Z,Y]$ для каждого горизонтального поля $Z$ ([3], 


с.\,59).


\end{pf}





{\bf 2.} В этом пункте вводится метрическая 


связность $\wt \nabla$, относительно которой


распределения $TF$ и $H$ являются параллельными, и изучаются 


проекции гладкой кривой $\gamma :I\to M$ в 


$L(p_0)$ и $S(p_0)$, где 


$p_0=\gamma (0)$, которые определяются 


с помощью связности $\wt \nabla$.


Так как распределение $H$ является связностью Эресмана 


для слоения $F$, существует единственная вертикально-горизонтальная 


гомотопия $P_\gamma :I\times I\to M$ такая, что 


$\gamma (t) =P_\gamma (t,t)$ для $t\in I$ [5].


В доказательстве теоремы де~Рама проекции кривой 


$\gamma$ в $L(p_0)$ и в $S(p_0)$ определены с помощью 


$\nabla$, и показано, что проекции 


$\gamma$ в $L(p_0)$ и в $S(p_0)$ совпадают с кривыми 


$P_\gamma (\cdot ,0):I\to M$ и 


$P_\gamma (0,\cdot ):I\to M$ соответственно. При этом используется


то, что распределение $H$ является вполне интегрируемым 


([1], с.\,180--183).


Ниже будет доказан аналогичный факт для проекций $\gamma$, 


определенных с помощью $\wt \nabla$ без предположения 


о том, что $H$ вполне интегрируемо. Доказательство дано 


для гладкой кривой $\gamma$, содержащейся в 


достаточно малой окрестности точки $p_0=\gamma (0)$ (теорема 1). 


В общем случае это утверждение легко получается 


из локального варианта делением отрезка $I$ на конечное число 


мелких частей $I_j$ таких, что для проекции 


кривой $\gamma :I_j\to M$ имеет место утверждение 


теоремы 1.


 


Теорема 2 утверждает, что связность $\wt \nabla$ совпадает с 


$\nabla$ тогда и только тогда, когда распределение $H$ 


вполне интегрируемо. 





Полагая $\wt \nabla_ZX=\nabla^1_ZX_1\oplus \nabla^2_ZX_2$, где $X$, 


$Z\in V(M)$, $X_i=\pi_i(X)$, $i=1, 2$, получим метрическую связность 


$\wt \nabla$ на $TM$. 


Нетрудно проверить, что распределения $TF$ и $H$ 


параллельны относительно $\wt \nabla$. 


Пусть $\gamma:I\to M$ --- гладкая кривая, $\gamma 


(0)=p_0$ и $\gamma (1)=p$, $C:I\to T_{p_0}M$ --- развертка кривой $\gamma$


в $T_{p_0}M$, определенная связностью $\wt \nabla$. (Cм. определение 


развертки в ([1], с.\,129). Здесь для удобства касательное векторное 


пространство $T_{p_0}M$ отождествляется с аффинным касательным 


пространством в точке $p_0$.) Пусть $C(t)=(A(t),B(t))$, где $A(t)\in F(p_0)$, 


$B(t)\in H(p_0)$ для $t\in I$. Так как $M$ --- полное риманово 


многообразие, $\wt \nabla$ --- метрическая связность, существуют гладкие 


кривые $\gamma_1, \gamma_2:I\to M$, которые развертываются на кривые 


$t\to A(t)$ и $t\to B(t)$ соответственно ([1], с.\,167, теорема 4.1). Согласно 


предложению 4.1 ([1], с.\,129) $\gamma_i$ --- такая кривая, что результат 


параллельного переноса $\Dot \gamma_i(t)$ в точку $p_0$ вдоль 


$\gamma_i^{-1}$, определенного связностью $\wt \nabla$, совпадает с 


результатом параллельного переноса $\pi_i(\Dot \gamma(t))$ в точку $p_0$ 


вдоль $\gamma^{-1}$, также определенного связностью $\wt \nabla$, где 


$i=1,2$. Следовательно, $\gamma_1$ является вертикальной кривой, а 


кривая $\gamma_2$ является горизонтальной. Кривые $\gamma_1$, 


$\gamma_2$ назовем проекциями кривой $\gamma$ в $L(p_0)$ и в 


$S(p_0)$ соответственно.





Пусть $q_0\in M$, $U$ --- достаточно малая связная относительно 


компактная окрестность точки $q_0$ в слое $L(q_0)$,


$B_\delta\to U$ --- расслоение, слой которого над $q\in U$есть


открытый шар в $H(q)$ радиуса


$\delta$. Тогда для достаточно малого $\delta >0$


экспоненциальные отображения связности $\nabla$, определенные в 


точках $U$, определяют диффеоморфизм


$$


\exp_U:B_\delta \to V


$$


множества $B_\delta$ на некоторую окрестность $V$ точки $q_0$ в $M$, 


содержащую $U$. Причем, если $p\in V$ и $p= \exp_q(X)$, $q \in U$, $X \in 


H(q)$, то расстояние $r(p,q)=|X|$ реализуется геодезической, чья 


скорость в точке $q$ равна $X$. При этом проекция 


$\pi^1:V\to U$, $\pi^1(p)=q$, является гладким тривиальным 


расслоением, слои которого диффеоморфны открытому шару в $R^{n-k}$. 


Обозначим через $Q(q)$ слой расслоения $\pi^1:V\to U$, проходящий 


через точку $q\in U$. Для достаточно малого $\delta >0$ слои слоения 


$F|_V$ и слои расслоения $\pi^1:V\to U$ пересекаются трансверсально. 


Следовательно, слоение $F|_V$ можно задать гладкой субмерсией 


$\pi^2:V\to Q(q_0)$, при которой $\pi^2(p)=q$, если $p\in L_q$, где $L_q$ --- 


слой слоения $F \mid_V$, проходящий через точку $q\in Q(q_0)$. Пусть 


$\gamma:I\to V$ --- гладкая кривая, $q_0=\gamma (0)$, $\mu=\pi^2(\gamma)$. В 


п.\,1 отмечалось, что распределение $H$ является связностью Эресмана 


для расслоения $\pi^2:V\to Q(q_0)$. Поэтому для каждого $t\in I$ 


существует единственное горизонтальное поднятие $P_t:I\to V$ с 


начальным условием $P_t (t)=\gamma (t)$ кривой $\mu :I\to Q(q_0)$. 


Полагая $P(t,s)=P_t (s)$, получим вертикально-горизонтальную гомотопию 


$P:I\times I\to M$ такую, что $\gamma (t)=P(t,t)$. Эта гомотопия 


единственным образом определяется кривой $\gamma$.





\begin{thm}


Проекция кривой $\gamma:I\to V$ в $L(q_0)$ $($в $S(q_0))$ есть 


кривая $P(\cdot,0):I\to L(q_0)$ $($соответственно $P(0,\cdot ):I\to S(q_0))$.


\end{thm}





\begin{pf}


Пусть $\Dot \gamma(t)=\Dot \gamma_1(t)\oplus \Dot 


\gamma_2(t)$ --- касательный вектор кривой $\gamma$ в точке $\gamma (t)$, 


где $\Dot\gamma_1(t)\in F(\gamma (t))$, $\Dot \gamma_2(t) \in H(\gamma (t))$, 


$t\in I$. В дальнейшем всюду под параллельным переносом касательного 


вектора будем понимать параллельный перенос, определенный связностью 


$\wt \nabla$. Для того чтобы показать, что проекция $\gamma$ в 


$L(q_0)$ (в $S(q_0)$) совпадает с $\nu= P(\cdot,0):I\to L(q_0)$ 


(соответственно с $h=P(0,\cdot):I\to S(q_0)$), достаточно показать, что 


параллельный перенос касательного вектора $\Dot \gamma_1(t)$ в точку 


$q_0$ вдоль $\gamma^{-1}$ совпадает с результатом параллельного 


переноса касательного вектора $\Dot \nu(t)$ в точку $q_0$ вдоль $\nu^{-


1}$ (соответственно параллельный перенос вектора $\Dot \gamma_2(t)$ в 


точку $q_0$ вдоль $\gamma^{-1}$ совпадает с параллельным переносом 


$\Dot h(t)$ в точку $q_0$ вдоль $h^{-1}$). Докажем этот факт сначала для 


проекции $\gamma$ в $S(q_0)$. Пусть $Z(t)$ --- горизонтальное поле


вдоль $\gamma:[0,\tau]\to M$, которое задает параллельный перенос


$\Dot\gamma_2(\tau)$ из $\gamma(\tau)$ в $q_0$. Будем считать $\tau=1$, что


не ограничивает общности. Тогда $\nabla^2_{\Dot \gamma}Z=0$ и $Z(1)=\Dot 


\gamma_2(1)$. Положим $Y(t,s)=D_tP(t,s)$, $X(t,s)=D_sP(t,s)$ для $(t,s)\in 


I\times I$. Очевидно, $\Dot \gamma(t)=Y(t,t)\oplus X(t,t)$ и 


$\nabla^2_{\Dot \gamma}Z=\pi_2[Y,Z]+\pi_2(\nabla_xZ)$ вдоль $\gamma$. 


Определим $Z$ во всех точках $P(s,t)$ следующим образом: вектор 


$Z(t)$ параллельно переносим из точки $\gamma (t)$ в точку $P(s,t)$ вдоль 


интегральной кривой векторного поля $Y$, проходящий через точку 


$\gamma(t)$, $t,s\in I$. Так как $\nabla^2$ --- метрическая связность, 


имеем $Yg(Z,Z)=0$. С другой стороны, в силу того, что $F$ риманово и 


$Z$ горизонтально, из ([3], с.\,47--61) имеем


$Yg(Z,Z)= 2g(Z,[Y,Z])$. Отсюда вытекает, что $[Y,Z]$ является 


вертикальным полем. Следовательно, $\nabla^2_{\Dot 


\gamma}Z=\pi_2(\nabla_xZ)$ и $\nabla^2_{\Dot \gamma}Z=0$ вдоль $\gamma$ 


влечет, что $\nabla_xZ$ является вертикальным полем вдоль $\gamma$. 





Положим $L^\prime(s)=(\pi^2)^{-1}(q)$ для $s \in I$, где $q=\mu(s)$. Тогда 


$L^\prime(q)$ является слоем расслоения $\pi^2:V\to Q(q_0)$ над точкой 


$q$. Семейство $\{L^\prime(s):s\in I\}$ образует слоение $F^\prime$ 


коразмерности один на многообразие $N^\prime=(\pi^2)^{-


1}(\Gamma_\mu)$, где $\Gamma_\mu$ --- образ кривой $\mu$, т.\,е. 


$\Gamma_\mu=\mu([0,1])$. Каждый слой слоения $F^\prime$ пересекает 


горизонтальную кривую $h=P(0, \cdot):I\to N^\prime$ трансверсально в 


одной точке. Поэтому слоение $F^\prime$ можно задать субмерсией 


$\pi:N^\prime\to \Gamma_h$, где $\Gamma_h$ --- образ кривой $h$. Пусть 


$f_{s_0}$ --- голономное отображение слоения $F^\prime$ вдоль кривой 


$v_{s_0}^{-1}$, где $v_{s_0}=P(\cdot, 


s_0):[0, t]\to L^\prime(s)$, $t\in I$. Тогда 


по определению голономного отображения имеем $f_{s_0}(P(t, s))=P(0, s)$ 


для $s$ близких к $s_0$. В силу того, что слоение $F$ риманово, слоение 


$F^\prime$ также является римановым.


Поэтому дифференциал $f^*_{s_0}$ отображения $f_{s_0}$ в точке 


$v_{s_0}(t)$ совпадает с параллельным переносом вдоль $v_{s_0}^{-1}$,


определенным связностью $\nabla^2$ [7].





Отсюда $\pi_*Z(t, s)=Z(0, s)$ для $t, s\in I$, где $\pi_*$ ---


дифференциал отображения $\pi$. По лемме 2 касательный вектор получен


параллельным 


переносом касательного вектора $\Dot \gamma_2(s)$ вдоль кривой $P(\cdot, 


s):[0, s]\to L^\prime(s)$. Поэтому имеет место $\Dot h(s)=\pi_*\Dot 


\gamma_2(s)$ для $s\in I$, и в частности, $Z(0, 1)=\Dot h(1)$. Следовательно, 


$\nabla_{\Dot h}Z=\nabla_{\pi_*\Dot \gamma_2}\pi_*Z$. В силу того, что 


$F^\prime $ --- риманово слоение, голономное отображение $f_{s_0}$ 


сохраняет длину горизонтальных кривых [8]. Поэтому по утверждению 3.1 


работы [8] имеем $\nabla_{\Dot h}Z(0, s) = \pi_* (\nabla_xZ(t,s))$. Так как 


$\nabla_xZ$ вертикально, $\pi_*(\nabla_xZ )=0$, т.\,е. $\nabla_{\Dot h}Z=0$, 


что в свою очередь влечет $\nabla^2_{\Dot h}Z=0$. 





Теперь докажем, что 


проекция $\gamma$ в $L(q_0)$ совпадает с $\nu$. Для удобства, как и 


выше, предположим $\tau=1$. Пусть параллельный перенос $\Dot 


\gamma_1(1)$ в точку $q_0$ вдоль $\gamma^{-1}$ задается вертикальным 


полем $Z$, определенным вдоль $\gamma$. Тогда 


$\nabla^1_{\Dot \gamma }Z=0$. Определим $Z$ во всех точках $P(t,s)$ следующим образом: 


вектор $Z(t)$ из точки $\gamma (t)$ параллельно переносим в точку 


$P(t,s)$ вдоль интегральной кривой векторного поля $X$, проходящий 


через точку $\gamma (t)$, $t,s\in I$. 





Пусть $X^*$ --- поле касательных векторов кривой $\mu$ в $Q(q_0)$. Не 


ограничивая общности, будем считать, что $X^*$ определено на $Q(q_0)$, 


т.\,к. всегда $X^*$ можем продолжить на некоторую окрестность кривой 


$\mu$. Поскольку распределение $H$ является связностью 


Эресмана для слоения $F$, векторное поле $X^*$ имеет горизонтальное 


поднятие на $V$, которое обозначим через $X$ (т.\,к. его сужение на 


поверхность $P:I \times I\to V$ совпадает с $X$). Таким образом, теперь 


векторное поле $X$ определено на $V$ и является горизонтальным 


поднятием векторного поля $X^*$. Поэтому каждый слой расслоения 


$\pi^2:V\to Q(q_0)$ под действием потока $X$ переходит в слой этого 


расслоения, следовательно, для каждого вертикального поля $W$ скобка 


Ли $[X, W]$ является вертикальным полем. Теперь вертикальное поле $Z$ 


продолжим на некоторую окрестность поверхности $P:I \times I \to N$. Из 


равенства $g([X, Z], W)=g(\nabla_XZ, W)-g(\nabla_ZX, W)$ для 


произвольного гладкого вертикального поля $W$ с учетом соотношений 


$g(\nabla_ZX, W)=Zg(X, W) - g(X, \nabla_ZW)$, $g(X, W)=g(X, 


\nabla_ZW)=0$ получаем $g([X, Z], W)=g(\nabla_XZ, W)$. Так как $W$ 


--- произвольное вертикальное поле, отсюда вытекает, что 


$\pi_1(\nabla_XZ)=\pi_1[X, Z]$. Так как $\nabla^1$ --- метрическая связность, 


имеет место $Xg(Z, Z)=0$ в точках поверхности $P:I\times I \to V$. С 


другой стороны $Xg(Z, Z)=2g(Z, \nabla_XZ)$, следовательно, в точках $P(t, 


s)$ векторные поля $\nabla_XZ$ и $[X, Z]$ являются горизонтальными полями. 


Отсюда вытекает, что $\pi_1(\nabla_XZ)=0$, и $[X, Z]=0$ в точках 


поверхности $P:I \times I\to V$. Значит, если учесть $\Dot \gamma 


=Y\oplus X$, то имеем $\pi_1(\nabla_{\Dot \gamma}Z)=\pi_1(\nabla_Y Z)$ 


вдоль $\gamma$. В силу вполне геодезичности $F$ векторное поле 


$\nabla_YZ$ является вертикальным, следовательно, $\nabla_YZ=0$ 


вдоль $\gamma$. Необходимо доказать, что $\nabla_{\Dot\nu} Z=0$. Пусть 


$f_s$ --- горизонтальное голономное отображение вдоль горизонтальной


кривой $h^{-1}_s$, 


где $h_s$ --- сужение $h$ на $[0, s]$, $s\in I$. По определению 


горизонтального голономного отображения, если $x$ --- точка из $(\pi^2)^{-


1}(q)$, $q=\mu(s)$, то точка $f_s(x)$ из $(\pi^2)^{-1}(q_0)$ определяется 


следующим образом: если $v_x:[0, 1]\to V$ --- гладкая вертикальная кривая, 


соединяющая точку $h(s)$ с точкой $x$, $Q:[0, 1]\times[0, s]\to V$ --- 


вертикально-горизонтальная гомотопия для пары $(v_x, h^{-1}_s)$, то 


$f_s(x) = Q(1, s)$. Известно, что $f_s$ является диффеоморфизмом слоев 


$(\pi^2)^{-1}(\mu(s))$ и $(\pi^2)^{-1}(q_0)$ расслоения $\pi^2:V\to Q(q_0)$. 


Для вполне геодезических слоений $f_s$ является изометрией [9]. 


Поэтому, если $(f_s)_*$ - дифференциал отображения $f_s$, то 


$(f_s)_*\nabla_YZ=\nabla_{(f_s)_*Y}(f_s)_*Z$ ([1], с.\,156).





Так как $[X, Z]=0$ в точках поверхности $P:I\times I\to M$, интегральная 


кривая векторного поля $Z$, начинающаяся в точках $P(t, s)$, под 


действием потока векторного поля $X$ переходит в интегральную кривую 


$Z$. Следовательно, $f_*Z(t, s)=Z(t, 0)$ для $(t, s)\in I \times I$. Также из 


$[X, Y]=0$ получаем


$(f_s)_*$ $Y(t,s)=Y(t,0) $ для $(t,s)\in I\times I$.


Отсюда вытекает


$(f_s)_*\nabla_YZ=\nabla_{\Dot v}Z$. Так как 


$\nabla_YZ=0$ вдоль $\gamma$, получим $\nabla_{\Dot v} Z=0$. Таким 


образом, $Z$ параллельно вдоль $v$. По лемме 2 вектор $\Dot v(1)$ 


получается параллельным переносом $Z(1, 1)$ вдоль интегральной кривой


векторного поля $X$. Поэтому $Z(1, 0)=\Dot v(1)$.


\end{pf}





\begin{thm}


Следующие утверждения эквивалентны.


\begin{itemize}


\item[1.] Распределение $H$ вполне интегрируемо


$($т.\,е. $\dim S(p)=n-k$ для каждого 


$p \in M)$.


\item[2.] $\wt \nabla$ является связностью без кручения $($т.\,е. $\wt 


\nabla=\nabla)$.


\end{itemize}


\end{thm}





{\bf Доказательство.} Покажем, что 1)$\Rightarrow$ 2). 





Пусть $H$ вполне интегрируемо. Тогда семейство $\{S(p), p \in M\}$ 


порождает $(n{-}k)$-мерное слоение $F^\bot$ многообразия $M$, которое 


является одновременно римановым и вполне геодезическим. Для любых 


векторных полей $X,Y\in V(H)$ имеем $[X,Y] \in V(H)$ (по теореме 


Фробениуса), $\nabla_XY \in V(H)$ (в силу того, что $F^\bot$ вполне 


геодезично) (см., напр., в [3]). Покажем, что $\wt \nabla$ является 


связностью без кручения, т.\,е. $\wt \nabla_XY-\wt \nabla_YX=[X,Y]$, 


для любых $Х, Y \in V(M)$. Вычислим левую часть требуемого равенства. 


Пусть $X=X_1\oplus X_2$, $Y=Y_1\oplus Y_2$, где $Х_i=\pi_iХ$, 


$Y_i=\pi_iY$, $i=1,2$. Тогда $\wt 


\nabla_XY=\nabla^1_XY_1+\nabla_X^2Y_2=\pi_1(\nabla_{X}Y_1)+\pi_2 


[X_1, Y_2)+ \pi_2 (\nabla_{X_2}Y_2) =\nabla_{X_1}Y_1+ 


\pi_1(\nabla_{X_2}Y_1) + \pi_2 [X_1, Y_2] + (\nabla_{X_2}Y_2)$.


Как отметили в п.\,1, $\pi_1(\nabla_{X_2}Y_1)=\pi_1[X_2, Y_1]$ из-за


горизонтальности $X_2$. Учитывая это, получим


$$


\wt \nabla_XY=\nabla_{X_1}Y_1+\pi_1[X_2, Y_1] +\nabla^2_{X}Y_2.


$$


Аналогично


$$


\wt \nabla_YX=\nabla_{Y_1}X_1+\pi_1[Y_2,X_1] + \nabla^2_{Y}X_2.


$$


Учитывая равенство $\nabla_XY-\nabla_YX=[X,Y]$, получим $\wt 


\nabla_XY - \wt \nabla_YX=[X,Y]$. Таким образом, $\wt \nabla$ ---


метрическая связность без кручения. Так как риманова метрика 


определяет единственную метрическую связность без кручения, имеем 


$\wt \nabla =\nabla$.





Теперь покажем, что 2)$\Rightarrow$1). Пусть $\wt \nabla =\nabla$. Тогда 


распределение $H$ параллельно относительно $\nabla$, следовательно, $H$ 


интегрируемо по предложению 5.1 из [1].


\medskip





{\bf Замечание.}


Как показывает известное расслоение Хопфа на трехмерной 


сфере, распределение $H$ не всегда вполне интегрируемо. В случае, когда 


$H$ вполне интегрируемо, имеет место теорема де~Рама: 


если $M$ односвязно, то $M$ изометрично произведению $L(p)\times S(p)$ 


для каждого $p\in M$. В этом случае $S(p)$ является слоем слоения 


$F^\bot$, порожденного распределением $H$. Определяются проекции 


произвольной точки $p\in M$ в $L(p_0)$ и в $S(p_0)$ для некоторого 


$p_0\in М$ следующим образом. Пусть $\gamma:I\to M$ --- гладкая кривая, 


$\gamma (0)= p_0$, $\gamma (1)= p$, $\nu$ --- проекция $\gamma$ в 


$L(p_0)$, $h$ --- проекция $\gamma$ в $S(p_0)$. Конечные точки кривых 


$\nu$ и $h$ назовем проекциями $р$ в $L(p_0)$ и в $S(p_0)$ 


соответственно. В силу того, что $H$ вполне интегрируемо, проекция 


$p$ зависит только от класса гомотопии кривой $\gamma$. Поэтому, когда 


$M$ односвязно, отображение $f: p\to (p_1, p_2)$ корректно определено. 


По теоремам 1 и 2 оно является изометрическим погружением. В силу


$\dim M=\dim\{L(p_0)\times S(p_0)\}$ отображение $f$ является накрытием,


следовательно, оно является изометрией ([1], с.\,134).








\begin{thebibliography}{9}





\bibitem{1}


Кобаяси Ш., Номидзу К. {\em Основы дифференциальной геометрии}. T.\,1. -- 


М.: Наука, 1981. -- 344~c.


\bibitem{2}


Reinhart B. {\em Foliated manifolds with bundle like metrics} // Ann. Math. 


-- 1959.


-- V.\,69. -- \No\,1. -- P.\,119--132.


\bibitem{3}


Tondeur Ph. {\em Foliations on Riemannian manifolds}. -- 


New York: Springer--Verlag, 1988. -- 247~p.


\bibitem{4}


Hermann R. {\em On the differential geometry of foliations} //


Ann. Math. Mech. -- 1960. -- V.\,72. -- \No\,3. -- P.\,445--457.


\bibitem{5}


Blumenthal R., Hebda J. {\em Ehresman connections for foliations} // Indiana 


Univ. Math. J. -- 1984. -- V.\,33. -- \No\,4. -- P.\,597--611.


\bibitem{6}


Blumenthal R., Hebda J. {\em Complementary distributions which preserve the


leaf geometry and applications to totally geodesic foliations} //


Quart. J. Math. -- 1984. -- V.\,35. -- P.\,383--392.


\bibitem{7}


Morgan A. {\em Holonomy and metric 


properties of foliations in higher codimension} //


Proc. Amer. Math. Soc. -- 1976. -- V.\,58. -- P.\,255--261.


\bibitem{8}


Hermann R. {\em A sufficient condition that a mapping of Riemannian manifolds


be a fibre bundle} // Proc. Amer. Math. Soc. -- 1960. -- V.\,11.


-- P.\,236--242.


\bibitem{9}


Johnson D., Whitt L. {\em Totally geodesic foliations on $3$-manifolds} //


Proc. Amer. Math. Soc. -- 1979. -- V.\,76. -- P.\,355--357.





\end{thebibliography}





\vskip 2cm





\postup{\it Ташкентский государственный}{$\qquad$\it Поступили}


\postup{\it университет}{{\it первый вариант} 17.12.1997}


\postup{}{{\it окончательный вариант} 29.10.1998}


\label{end}


\end{document}





